\section{Gricigliano and the Formation of a Common Culture}

Gricigliano, in the hills outside Florence, became the place where a dispersed founding network acquired a durable center. The property, associated with the Villa Martelli, had been used by Benedictines of Fontgombault. Institutional memory says the previous benefactors expected continued celebration of the older Mass, that Cardinal Augustin Mayer helped connect the parties, and that Archbishop Silvano Piovanelli received the Institute's motherhouse and seminary in the Archdiocese of Florence. The exact deed, conditions, and diocesan erection act were not located.

The Saint Philip Neri Seminary is not merely a school attached to an apostolic network. It gathers candidates from multiple countries into a common language, schedule, intellectual program, liturgical life, and government. Current institutional sources describe a multi-year philosophical and theological course rooted in Roman ecclesial culture. The 2026 calendar names Philippe Mora as seminary superior, maintaining the co-founder's role at the center of formation while Wach serves as prior general.

Formation explains how an institution can expand without becoming a federation of unrelated local initiatives. A candidate formed in Tuscany may later serve a restored urban shrine, a parish-like apostolate, a school, or an African mission. The same seminary culture joins liturgy, preaching, catechesis, priestly fraternity, and a cultivated attention to sacred art and music. That coherence is an institutional aspiration; the checked public record does not include formation evaluations or a survey of members showing uniform reception.

\subsection{The language of \emph{Romanità}}

The Institute uses \emph{Romanità} to name more than residence in Rome. Its public presentation connects fidelity to the pope, Catholic doctrine, Latin liturgy, the Church's councils and canons, and a style of culture and government. In this narrative the founding network around Roman cardinals, the founders' studies, and the seminary's proximity to Rome all reinforce one self-understanding.

Historical analysis must keep the term attributed. \emph{Romanità} is the Institute's synthesis, not a measurement proving that all its judgments are those of the Roman Curia. Roman approval in 2008 and 2016 establishes acts of authority; institutional language about Roman spirit establishes identity. Neither makes a private audience a decree or an apostolate immune from its diocesan bishop.

\subsection{Beauty, discipline, and pastoral purpose}

Gricigliano's visual and ceremonial culture became part of the Institute's public signature. Solemn liturgy, choir, vesture, architecture, processions, and a canonial style communicate continuity and hierarchy. The Institute argues that beauty disposes people toward divine things and integrates grace, truth, charity, and culture.

That claim can be recorded without turning aesthetics into proof of sanctity or effectiveness. A historical account asks how the program organized formation and reception. The public sources establish an intended coherence; without formation evaluations or a member survey, this study cannot measure uniform reception or assess individual practices. The later visitation and public scrutiny are separate oversight records, not verdicts for or against the Institute's aesthetic ideals.

\evidenceframe{ICKSP presentation pages and 2026 international calendar; institutional history of Gricigliano.}{The motherhouse and international seminary near Florence became the durable center of formation, government, and common culture; Mora is listed as seminary superior in 2026.}{The original property and erection instruments, full curriculum, internal evaluations, and formation directories were not located. Public self-description is not an external quality audit.}
